% Generated draft: requires mathematical and source review.
\documentclass[11pt]{article}
\usepackage[margin=25mm]{geometry}
\usepackage{fontspec}
\setmainfont{Latin Modern Roman}
\usepackage{amsmath,amssymb,mathtools,unicode-math}
\setmathfont{Latin Modern Math}
\usepackage{enumitem,needspace,xurl,hyperref}
\setlength{\parindent}{0pt}
\setlength{\parskip}{5pt}
\setlength{\emergencystretch}{3em}
\newcommand{\N}{\mathbb{N}}
\newcommand{\Z}{\mathbb{Z}}
\newcommand{\Q}{\mathbb{Q}}
\newcommand{\R}{\mathbb{R}}
\newcommand{\C}{\mathbb{C}}
\newcommand{\F}{\mathbb{F}}
\newcommand{\A}{\mathbb{A}}
\newcommand{\PP}{\mathbb{P}}
\newcommand{\E}{\mathbb{E}}
\newcommand{\eps}{\varepsilon}
\newcommand{\dd}{\,\mathrm{d}}
\begin{document}
\section*{UnsolvedMath ID: 30004846}
\textbf{Source problem code:} catalogue entry
\subsection{Definitions and mathematical statement}
A fatgraph is a graph equipped with a cyclic ordering of the half-edges incident to each vertex. The supplied source describes a \emph{short-branched tree} as a metric fatgraph tree satisfying a particular length condition, but does not give that condition. It describes a \emph{string diagram} as a metric fatgraph assembled from disjoint input circles and collections of short-branched trees attached inductively along their leaves. The resulting diagrams form a cell complex denoted by \(SD\). The excerpt does not specify the full metric conventions, permitted degenerations, cell identifications, or component parameters needed to define this complex independently.

Let \(\mathcal M_{\partial}\) denote the source's intended moduli space of Riemann surfaces with boundary. This notation leaves unresolved the topological types of the surfaces, whether surfaces must be connected, boundary labels and parametrizations, any division into incoming and outgoing boundaries, the equivalence relation on surfaces, and the topology on the moduli space. In particular, the record does not say whether the assertion concerns fixed components or unions of components.

\textbf{Conjecture as supplied.} For the intended choices of \(SD\) and \(\mathcal M_{\partial}\),
\[
SD\simeq\mathcal M_{\partial}.
\]
Explicitly, the proposed homotopy equivalence means that there exist continuous maps
\[
f\colon SD\longrightarrow\mathcal M_{\partial},\qquad
 g\colon\mathcal M_{\partial}\longrightarrow SD
\]
such that
\[
g\circ f\simeq\operatorname{id}_{SD},\qquad
f\circ g\simeq\operatorname{id}_{\mathcal M_{\partial}}.
\]
The missing conventions above must be recovered before this becomes a fully specified mathematical assertion; no choice of those conventions is made here.

The source motivates the conjecture by relating the short-branched-tree metric condition to an equivalence relation in Bodigheimer's work on a harmonic compactification of moduli space. It says that these trees interpolate, in an unspecified sense, between equivalent cells, and that compelling examples of face identifications of assocoipahedra produce small circles near the harmonic boundary. These are supplied motivations, not a definition of the target as that compactification.

The excerpt also reports joint work with G.C. Drummond-Cole and N. Rounds in which \(SD\) parametrizes operations on singular chains of the loop space of a closed, oriented manifold, with composition respected on homology. Separately, it describes work in progress with T. Tradler using assocoipahedra to construct a homotopy equivalence \(N\,DG\to SD\). Neither contextual assertion resolves the unspecified moduli-space conventions or establishes the conjecture here.
\subsection{Short English statement}
The supplied conjecture asks whether the string-diagram cell complex \(SD\) has the homotopy type of the intended moduli space of Riemann surfaces with boundary. Precise definitions of the diagram complex, surface types, boundary data, and moduli topology remain missing from the excerpt. This is a statement-only draft requiring human mathematical review.
\subsection{Sources}
\begin{itemize}
\item Supplied local UnsolvedMath registry record, permanent ID \(30004846\), title \emph{String Diagram Models of Moduli Spaces}, source problem code \texttt{catalogue entry}, in the Oberwolfach Reports Open Problems collection. The record supplies proposed year 2021 and no named proposer.
\item Supplied report citation: \emph{Homotopical Algebra and Higher Structures} (2022), Oberwolfach Reports, DOI 10.4171/owr/2021/46, pp.\ 2497--2560. Source URL and literature-source label ``Oberwolfach Report 46'': \url{https://doi.org/10.4171/owr/2021/46}.
\item Supplied literature source: \emph{Chain-level string topology operations}, \url{https://arxiv.org/abs/1506.02596}. The record does not explicitly identify this link with either numbered reference in the excerpt.
\item The original excerpt refers to Bodigheimer's work as reference [1] and to joint work with G.C. Drummond-Cole and N. Rounds as reference [2], but supplies neither numbered bibliography entry. It separately attributes work in progress to collaboration with T. Tradler; the speaker is not identified in the supplied record.
\item Literature assessment dated exactly \texttt{2026-08-22}: the supplied curation labels the problem open and states that string diagrams support chain-level string-topology operations and related graph models have been compared, while the stated homotopy equivalence remains conjectural. This is supplied dated curation derived from the OWR record, not an independently verified literature assessment. Its time-sensitive status and cited claims require independent verification.
\end{itemize}
\end{document}
